\section{A marked specialization gives a seed on the same determinant component}

Fix an original setup $T$ of rank $n=2$ or $3$ with coefficient DVR
$O$, residue field $k$ and frame $F$. Write $D_0=T.B.B$ for the
universal ordinary determinant ring, $A=T.S.A$ for the matrix entry
ring, and $B=T.S.B$ for its coefficient subring. Fix any comparison
$\Phi:D_0\to B$ satisfying $T.\mathrm{IsPhi}$. The twist link
has a continuous surjection $p_0:G_0\to F.G_0$, a character
$\kappa:G_0\to O^\times$, and local identifications
$\lambda_v:G_v\simeq F.G_{e(v)}$. Twisting a representation $b$
of $F.G_0$ gives $g\mapsto\kappa(g)b(p_0g)$. Use $g^*$ for
the inverse of conjugation applied to $g$, on either source group.
The two group operations commute with $p_0$.

For the ideal $J\subset D_0$ assume
\[
 T.\mathrm{Unr0}(J),\qquad T.\mathrm{DimT}(J),\qquad
 T.\mathrm{IdealSpecializationPoint}(J).
\]
The last condition gives an algebraically closed $O$-field $K'$,
a map $\varphi:D_0\to K'$ killing $J$, and an irreducible
normalized $\rho'$ on $F.G_0$. At $\varphi$ the universal
determinant is the determinant of $\kappa\cdot(\rho'\circ p_0)$.
Only this determinant identity is part of the given point.

% BEGIN S68D1
\paragraph{S68D1. A generic irreducible representation on a component through the point.}
Choose a prime $Z_0$ minimal over $J$ and contained in $\ker\varphi$.
Put $C=D_0/Z_0$, let $q:D_0\to C$ be the quotient and
$K=\overline{\operatorname{Frac}C}$. Thus $C$ is a complete
Noetherian local domain with residue field $k$, it embeds in $K$,
and $\varphi$ factors as $\varphi' q$. The universal determinant
and all its ordered characters specialize under $q$ to an ordinary
pair $(D,\psi)$ over $C$. Its residual, ordinary, polarization
and inertia identities follow by base change of the pair; its
continuity descends modulo every power of the maximal ideal.

Over $K$ the determinant law $D_K$ is that of a semisimple
$r_0:G_0\to\mathrm{GL}_n(K)$. This is the determinant
representation theorem over an algebraically closed field
(Chenevier, \emph{The $p$-adic analytic space of pseudocharacters
of a profinite group, and pseudorepresentations over arbitrary rings},
Theorem~A). It applies also in positive characteristic. In the
present case $r_0(G_0)$ spans $M_n(K)$. Indeed the twist of
$\rho'$ spans $M_n(K')$ by Burnside, the surjectivity of $p_0$,
and pointwise multiplication by nonzero scalars. Choose group
elements $h_1,\ldots,h_{n^2}$ giving a basis there. Matrix units
show that $(X,Y)\mapsto\operatorname{tr}(XY)$ is perfect in
every characteristic. Hence
\[
 \Delta=\det\bigl(\operatorname{tr}_{D}(h_i h_j)\bigr)_{i,j}
 \quad\hbox{has nonzero image under }\varphi'.
\]
It is nonzero in $C$ and remains nonzero in $K$. The same group
elements give a basis of $M_n(K)$ under $r_0$. In particular if
$\operatorname{tr}(r_0(h)X)=0$ for every $h$, then $X=0$.

For $g\in G.I_0$, the unramifiedness premise puts $g-1$ in
the determinant kernel over $D_0/J$ and hence after passage
to $C$ and $K$. The coefficient of the linear term in the
determinant identity gives
$\operatorname{tr}(r_0(h)(r_0(g)-1))=0$ for every $h$
(evaluate the kernel identity at $1+s h(g-1)$). Thus
$r_0(g)=1$. All these operations use the same $Z_0$.
% END S68D1

% BEGIN S68D2
\paragraph{S68D2. Descent through the prescribed frame quotient.}
Put $a(g)=\kappa(g)^{-1}r_0(g)$ and
$t_0(g)=\kappa(g)^{-1}\operatorname{tr}_{D}(g)\in C$.
The trace of $a$ is $t_0$ in $K$. Its group matrices still span
$M_n(K)$. The residual identity of the pair gives
$\overline{t_0(g)}=\sum_i\Psi(p_0g)^{w_i}$.
Polarization of $D$ and
$\kappa(g)\kappa(\operatorname{conj}g)=\varepsilon(g)^{1-n}$
give $t_0(g^*)=t_0(g)$. The trace and the twisting character
are continuous modulo each power of $\mathfrak m_C$.

Let $N=\ker a$. By trace perfection,
\[
 g\in N\quad\Longleftrightarrow\quad
 t_0(hg)=t_0(h)\quad\text{for every }h\in G_0.
\]
Thus $N$ is a closed normal subgroup: each trace fiber is closed
by adic continuity and separatedness of $C$. It is stable under
conjugation by the specified involution. Indeed star reverses
products and fixes $t_0$, and cyclic invariance gives
$t_0(hg^*)=t_0(g h^*)=t_0(h^*g)=t_0(h^*)$ for $g\in N$.
Since $N$ contains inverses this proves the assertion for
conjugation as well. It contains $G.I_0$ by S68D1 and the
link identity $\kappa|_{G.I_0}=1$.

For each tame place $u$ and $g\in G.I_u$ with $\Psi(p_0g)=1$,
the link supplies $m$ prime to $\ell$ with $(p_0g)^m=1$ and
$\kappa(g)^m=1$. Hence $g^m$ lies in the prescribed inertia
kernel of every residual character. The pair's inertia equation
gives $\operatorname{tr}(r_0(h)(r_0(g^m)-1))=0$ for all $h$,
so $a(g)^m=1$. The polynomial supplied by $D$ on
$\kappa(g)^{-1}[g]$ is its monic characteristic polynomial over
$K$ and reduces to $(X-1)^n$ in $k[X]$. The polynomial
$Q=(X^m-1)/(X-1)$ has $Q(1)=m\in k^\times$. Consequently its
resultant with this characteristic polynomial is a unit of $C$.
B\'ezout and Cayley--Hamilton make $Q(a(g))$ invertible, and
$(a(g)-1)Q(a(g))=0$ gives $a(g)=1$. The original kernel axiom
of the twist link now gives $\ker p_0\subseteq N$.
There is therefore a representation $b:F.G_0\to\mathrm{GL}_n(K)$
with $b(p_0g)=a(g)$. Its trace descends to $t:F.G_0\to C$ by
the injectivity $C\hookrightarrow K$. The group projection $p_0$
is a quotient map (continuous surjection of compact Hausdorff
groups), so $t$ is adically continuous. It has the stated residual
trace, is star-invariant, and $\varphi'(t(x))=\operatorname{tr}\rho'(x)$
for all $x$, by the specialization identity in the hypothesis.
% END S68D2

% BEGIN S68D3
\paragraph{S68D3. Recovering the ordinary matrices and their symmetric polarization.}
In the local identification at $v$, define
\[
 \chi_{e(v),i}(\lambda_v s)=\psi_{v,i}(s)\,
                         \kappa(\phi_v s)^{-1},
\]
where $\phi_v:G_v\to G_0$ is the original local map.
These $C$-valued unit characters are adically continuous, reduce to
$\Psi(\operatorname{loc}_{e(v)}(-))^{w_i}$, and give the original
$\psi_{v,i}$ on retwisting. The ordinary equation of the pair
annihilates the trace of each group translate of its ordered
ordinary product. S68D1 shows that the matrix product itself
vanishes under $r_0$. Remove the scalar unit $\kappa(\phi_v s_i)$
from the $i$th factor. This leaves the required product for $b$
and $\chi$ in precisely its original order and for independently
chosen arguments $s_i$.

The representation $g\mapsto b(g^*)^t$ and $b$ have equal
traces, and each spans the full matrix algebra. Their two maps
from $K[F.G_0]$ have the same kernel: if one kills $y$, evaluate
the common trace on $hy$ for all $h$ and use trace perfection
for the other. The resulting automorphism of $M_n(K)$ is inner by
Skolem--Noether, giving an invertible matrix $P$ such that
$b(g^*)^tP=P b(g)$. Its transpose also intertwines by involutivity,
so $P^t=\pm P$ by the scalar commutant and transposing twice.
Here $2$ is a unit of $C$, hence of $K$ and $K'$, as $F.\mathrm{Ok}$
requires odd residue characteristic. In rank three an invertible
alternating form is impossible. In rank two the alternating
alternative would give
$b(g)+b(g^*)=\operatorname{tr}(b(g))I$ by the formula for the
adjoint of a two-by-two matrix in an alternating form. Take traces
after multiplication by $b(x)$. The identity in $C$ specializes
at $\varphi'$ to the corresponding identity for $\rho'$, using
S68D2. Its full matrix span recovers the matrix identity at $\rho'$.
The given normalized point has
$\rho'(g^*)=J_n\rho'(g)^tJ_n$ for the standard anti-diagonal
$J_n$. In rank two, the off-diagonal entries of
$\rho'(g)+J_2\rho'(g)^tJ_2$ are twice those of $\rho'(g)$.
Thus all $\rho'(g)$ would be diagonal, contradicting irreducibility.
The generic form is symmetric in both ranks.
% END S68D3

% BEGIN S68D4
\paragraph{S68D4. Normalizing the marked difference over the same component.}
Set $D_* = b(g_0)-b(g_0^*)$, $z=\Psi(g_0)-\Psi(g_0)^{-1}$
and $X=t(g_0^2)-t(g_0g_0^*)$. The separator gives $z\ne0$;
using the weights $(1,-1)$ or $(1,0,-1)$ shows $\bar X=z^2$.
Hensel's lemma in $C$ gives $\zeta^2=X$ with $\bar\zeta=z$.
The difference $D_*$ is skew-adjoint for $P$, and
$\operatorname{tr}(D_*^2)=2X$ by star invariance and cyclicity.
Consequently its characteristic polynomial is $Y^2-X$ in rank
two and $Y(Y^2-X)$ in rank three. Indeed skew-adjointness pairs
the characteristic roots under sign, and the coefficient of
$Y^{n-2}$ is $-\operatorname{tr}(D_*^2)/2$. This works in
residue characteristic three as well.

Order the distinct eigenlines by the eigenvalues $w_i\zeta$.
The symmetric pairing vanishes on lines whose eigenvalues do
not sum to zero. Opposite lines therefore pair nontrivially;
scale their vectors to pairing one. In rank three the middle line
has nonzero norm, which can be scaled to one over the algebraically
closed $K$. In this basis the form is $J_n$ and $D_*$ is
$\operatorname{diag}(w_i\zeta)$. With the characters $\chi$ of
S68D3 the conjugated $b$ defines an irreducible normalized
representation $\rho$ of the original frame. Its trace, ordered
characters and marked difference satisfy all three original
$\mathrm{Integral}(O,C,\rho)$ conditions. Its twist is conjugate
to $r_0$. Hence its determinant is the image of $D$, and its
twisted characters are exactly the images of $\psi$.
% END S68D4

% BEGIN S68D5
\paragraph{S68D5. Realizing the integral generic point on the original entry algebra.}
Let $\mathcal E$ be the image of $C[F.G_0]$ under $\rho$.
Every element has trace in $C$, and the reduction of this trace
is obtained by evaluating the same group-algebra word on the
diagonal residual representation with characters $\Psi^{w_i}$.
Indeed this holds on group elements by S68D2, hence on all
finite $C$-linear combinations. The distinct elements $w_i\zeta$
have pairwise unit differences in $C$. The interpolation
polynomials for $D_*$ therefore give the diagonal matrix units
$e_i\in\mathcal E$. On the diagonal residual representation the
same words give the corresponding coordinate projectors:
the residual marked difference has entries $w_i z$.

These projectors make the cone conditions explicit. First,
$\rho(g)_{ii}=\operatorname{tr}(e_i\rho(g))$ belongs to $C$
with the correct residue. Any product around a directed cycle
of group entries is the trace of the word obtained by inserting
the coordinate projectors between its factors. Its value lies
in $C$ and reduces to zero if the cycle traverses an off-diagonal
edge, since the residual representation is diagonal. Write
$N_{ij}$ for the $C$-span of the entries in position $(i,j)$.
Polarization identifies the entries $(i,j)$ with those in
$(j^{\rm rev},i^{\rm rev})$ evaluated on the starred group element.
For rank two every weight-zero monomial with an off-diagonal
factor contains pairs of opposite off-diagonal factors. In rank
three remove the diagonal factors. Each list of nonzero degrees
summing to zero decomposes into pairs $(1,-1),(2,-2)$ and triples
$(1,1,-2),(-1,-1,2)$: first cancel equal absolute values, then
group the remaining ones in pairs. In each part the factors can
be arranged in a directed cycle. For the edges of degree one
one can use either $01$ or $12$ by polarization; for degree minus
one use either $10$ or $21$. The outer edges have unique positions.
Thus every balanced monomial with an off-diagonal factor is in
$\mathfrak m_C$. The all-diagonal factors already belong to $C$.

For each $i,j$ choose a nonzero $y=\rho(h)_{ji}$, possible since
the matrices span $M_n(K)$. Multiplication by $y$ embeds $N_{ij}$
as an ideal $U\subset C$, by the two-edge trace computation
(or the diagonal computation if $i=j$). This proves finite
generation. The function $g\mapsto y\rho(g)_{ij}$ is adically
continuous in $C$: it is $\operatorname{tr}(e_j\rho(h)e_i\rho(g))$,
a finite $C$-linear combination of translates of $t$. For every
$r$, Artin--Rees gives $s$ with
$U\cap\mathfrak m_C^s\subseteq\mathfrak m_C^r U$.
Cancelling $y$ proves local constancy in the required quotient
$N_{ij}/\mathfrak m_C^rN_{ij}$. The remaining character
continuity and residues were proved in S68D3.
Thus all five original cone conditions hold over $C\hookrightarrow K$.
APD10 gives $O$-maps $\xi:A\to K$ and $\beta:B\to C$ with
$\xi_*u=\rho$ and $\xi|_B=\beta$ inside $K$.
% END S68D5

% BEGIN S68D6
\paragraph{S68D6. The two seed primes and their dimension.}
We claim that $\beta\Phi=q$ on the entire determinant ring.
The comparison identity for $\Phi$ lives initially in the complete
universal matrix ring. Both sides come from the entry subring
$A$, so it holds already there: characteristic-law coefficients
agree under the injection of $A$, and so do the values of every
ordered character. Base change by $\xi$ now gives the determinant
and the twisted characters of $\rho$ at $\beta\Phi$ in $K$.
By S68D4 these are also the specialization under $q$ of the
universal pair. Since $C\hookrightarrow K$ is injective, equality
holds over $C$ (for the laws, compare all characteristic
coefficients on words). Universality for the pair over the
complete coefficient ring $C$ now gives $\beta\Phi=q$.
In particular $\beta$ is surjective, even though no surjectivity
was assumed for $\Phi$.

Take $\mathfrak p=\ker\xi\subset A$ and $Z=\ker\beta\subset B$.
They are prime, $Z=\mathfrak p\cap B$, and
$J B\subseteq Z$ where $JB$ denotes the ideal generated by
$\Phi(J)$. For every position $(i,j)$ some entry survives at
$\mathfrak p$, since $\rho(F.G_0)$ spans the matrix algebra.
Thus both triangular ideals survive, and in rank three an adjacent
entry survives as well. This is exactly type I in both ranks.
Lastly $B/Z\simeq C=D_0/Z_0$. The hypothesis $T.\mathrm{DimT}(J)$
applied to this minimal prime gives
\[
                 \dim B/Z=1+nF.d .
\]
These data prove the full
$\mathrm{SeedIn}(T.S,J B)$, for every original $T$ and every
comparison $\Phi$ satisfying $T.\mathrm{IsPhi}$.
This is $T.\mathrm{CurveSpecializationHypothesis}$.

In the linked arithmetic construction the specialization premise
has its own inputs. The native $\mathrm{SeedPtOfAlt}$ takes a
concrete realization $R$ with its specified character $\kappa_a$,
$\mathrm{IsKa}$, a prime-to-$\ell$ order character $\chi'$ and
the required cyclotomic congruence at every element, a spacing
$m\ge2$ and the matching $\mathrm{IsSeedAlt}$ representation.
With global Artin reciprocity it turns
$R.B_0.\mathrm{PointTheoremBaseObject}(\iota,J)$ into
$R.T.\mathrm{IdealSpecializationPoint}(J)$ for the very same $J$.
The present result uses that specialization together with
$\mathrm{Unr0}$ and $\mathrm{DimT}$.
% END S68D6
